\section{Research Objectives}

\subsection{General Objective}
The primary objective of this research is to design, implement, and evaluate a robust, non-parametric, retrieval-augmented continual learning framework for Vietnamese legal question-answering and reasoning. By maintaining a completely frozen base Large Language Model and updating an external knowledge store, the framework aims to enable real-time legal knowledge updates and controllable unlearning of confidential or repealed laws without suffering from catastrophic forgetting or incurring parameter retraining overheads.

\subsection{Specific Objectives}
To achieve the general objective, the project is structured around three specific research goals:
\begin{enumerate}
    \item \textbf{Develop a Selective RAG Memory System (Objective 1):} Formulate a dynamic legal importance score integrating document hierarchy, citation network centrality, and legal validity. Implement an Ebbinghaus-based consolidation and decay mechanism to manage storage tiers and reduce retrieval noise, a compliance hard-gate to guarantee absolute coverage of active laws, and the temporal/compliance gating---including the Retriever Policy ($P$) and Prompt Constraint ($Q$) unlearning filters for replaced legislation, sealed cases, and personal data---that keeps retrieval grounded on in-force law.
    \item \textbf{Quantify Catastrophic Forgetting under Continual Ingestion (Objective 2):} Build parametric comparison baselines---Continual Pre-Training / Continual Instruction Tuning and the semi-parametric Retrievable Gradients (ReGrad) approach \parencite{su2026retrievablegradientscontinualposttraining}---at the same small model scale, ingest the legal corpus in incremental snapshots, and measure the catastrophic forgetting (Backward Transfer and performance drift across ingestion volume) of each paradigm against the frozen RAG framework on QA, NLI, and syllogism.
    \item \textbf{Conduct Standardized Benchmarking and Validation (Objective 3):} Establish quantitative evaluation pipelines using Vietnamese legal datasets, namely VLQA \parencite{nguyen2025vlqacomprehensivelargehighquality}, ViLegalNLI \parencite{duong2026vilegalnlinaturallanguageinference}, and LegalSLM \parencite{le-etal-2025-overview}, with concrete test scenarios sourced from data-rich legal sub-domains such as labour and tax law, to validate the framework's retrieval accuracy, general utility retention, forgetting under sequential updates, and computational cost-efficiency, and to characterize the task-dependent reasoning ceiling along the QA $\rightarrow$ NLI $\rightarrow$ syllogism progression.
\end{enumerate}
